\section{Methodology}\label{sec:methodology}

\subsection{What the reported times contain}\label{subsec:timing}

Every phase is timed on the host with \texttt{time.perf\_counter} and recorded per run; Table~\ref{tab:timing} lists
the phases and states which of them each reported time includes. Device open and bitstream load ($T_{\mathrm{setup}}$)
happen once per session and are reported separately. $T_{\mathrm{e2e}}$ is the time from the Qiskit circuit object to
the sampled measurement results (1024 shots) and includes compilation, all transfers, all kernel time, all exchanges,
the final read back and sampling; it excludes $T_{\mathrm{setup}}$ and the reference computation used for fidelity.
For the baseline design the same boundaries are used; its worker processes are started once per session.

\begin{table}[t]
\caption{Timing phases and their inclusion in the reported times.}\label{tab:timing}
% generated by scripts/make_tables.py, do not edit
\begin{tabular}{@{}llcc@{}}
\toprule
Phase & Contents & in $T_{\mathrm{e2e}}$ & in $T_{\mathrm{kernel}}$ \\
\midrule
$T_{\mathrm{setup}}$ & device open, bitstream load, worker start (once per session) & no & no \\
$T_{\mathrm{compile}}$ & Qiskit circuit to engine records (both F4 schedules) & yes & no \\
$T_{\mathrm{h2d}}$ & initial state on the device, instruction uploads & yes & no \\
$T_{\mathrm{kernel}}$ & all engine launches, start to completion of the last engine & yes & yes \\
$T_{\mathrm{exch}}$ & remaps: device copies, host exchanges (PCIe and host copy), relabelling & yes & no \\
$T_{\mathrm{d2h}}$ & final read back of all slices & yes & no \\
$T_{\mathrm{post}}$ & sampling of 1024 shots from the slices & yes & no \\
reference & Qiskit or Aer reference state and error metrics & no & no \\
\bottomrule
\end{tabular}

\end{table}

\subsection{Protocol}\label{subsec:protocol}

Each configuration runs one warm up repetition, which is discarded, followed by \NumRepsMin\ to \NumRepsMax~timed
repetitions (more for short runs, fewer for the widest runs, whose length is bounded by the thermal limit of
Section~\ref{sec:limitations}). We report the median and the interquartile range.
Runs were made with no other job on the host (no build, no other benchmark; checked and logged before every campaign).
The CPU frequency governor was \NumGovernor\ and is recorded with every campaign. All circuits come from one
generator with fixed seeds (Section~\ref{subsec:compilation}); random input states are drawn from the Haar measure with
a fixed seed per width.

\subsection{Correctness metrics}\label{subsec:metrics}

For every reported instance the final state $\psi$ (read back from HBM, outside the timed region) is compared with
an independent reference $\psi_{\mathrm{ref}}$: the Qiskit \texttt{Statevector} class for $n \le \NumRefSvMax$ and Qiskit
Aer in double precision with gate fusion disabled above. We report the fidelity
$F = |\langle \psi_{\mathrm{ref}} | \psi \rangle|^2$, the infidelity computed without cancellation as
$1 - F = d^2 - d^4/4$ with $d = \min_\phi \lVert \psi - e^{i\phi} \psi_{\mathrm{ref}} \rVert$ for normalised states, the
largest amplitude error after optimal global phase alignment, and the $\ell_2$ error $d$. A global phase is
physically irrelevant and is removed by the alignment; it is nonetheless reported, since the baseline design's gate
decompositions drop global phases. Grover runs additionally report the measured target probability against the
analytic value $\sin^2((2k+1)\theta)$, and the phase estimation runs the full distribution of the counting register.

\subsection{CPU reference timing}\label{subsec:method-cpu}

To compare the FPGA simulation time with the Qiskit simulation time, on the same host and
for exactly the same circuit objects we time (i) the Qiskit \texttt{Statevector} class (NumPy, complex128, one
thread) and (ii) Qiskit Aer, statevector method, double precision, with its default thread count and with one
thread (versions in Table~\ref{tab:platform}). For Aer the circuit is transpiled to the simulator's native gate set
with \texttt{optimization\_level=0}, which keeps multi controlled gates such as the Grover oracle as single gates, and
Aer's default gate fusion is left enabled, as a user would run it; the reference states for correctness are computed
with fusion disabled. The FPGA engine receives the same circuit object and compiles a multi controlled gate into
one record (Section~\ref{subsec:compilation}); the baseline design lowers it to its 19 opcodes with the decomposition
of its original code. Only the simulation call is timed; circuit construction and transpilation
are excluded for both CPU and FPGA, and the FPGA time we compare against is $T_{\mathrm{e2e}}$, which does include the
FPGA side compilation. A comparison with GPU simulators is not repeated here: a detailed FPGA, CPU and CUDA-Q GPU
comparison on the same type of card is the subject of the companion single card study~\cite{tiwari2026fpga30}, and this
paper makes no speed claim against GPU simulators.

\subsection{Energy}\label{subsec:method-energy}

Board power of each card is read from its management controller (sum of the 12~V PCIe, 12~V auxiliary and 3.3~V rails,
the quantity reported by \texttt{xbutil}) every 0.1~s during a dedicated energy stage, in which every configuration is
repeated several times. The controller refreshes its readings about once per second. For every run the board power of a
card is the time weighted mean of the readings in force during the run, including the reading taken before the run
started, and the energy is this power times the end to end time, summed over the cards in use; we report the median
over the repetitions together with the number of new readings that fell inside the shortest run (Table~\ref{tab:energy}).
Host CPU and memory power are not measured, because the RAPL energy counters of this host are readable only with
administrator rights, so the energy figures cover the accelerator boards only.
